\section{Introduction}
\label{sec:intro}

\subsection{Motivation}
\label{sec:intro:motivation}

Recent advancements in large language models extend context length through hybrid architectures. Models like Jamba combine transformer attention mechanisms with State Space Models (SSMs) to support a 256K context window \cite{lieber2024jamba}. The base Jamba architecture also uses a Mixture-of-Experts (MoE) feedforward configuration, though our work focuses specifically on memory management for the attention and SSM cache types; MoE routing memory is orthogonal to this design. This architecture merges two different layer types: attention mechanisms requiring a linearly scaling Key-Value (KV) cache ($O(n)$ per token), and SSMs requiring a fixed-size state footprint ($O(1)$ per layer) \cite{dao2024mamba2}.

Existing inference engines struggle to serve these heterogeneous memory profiles. The unified pool approach in vLLM \cite{kwon2023pagedattention} forces the fixed SSM state to pad up to the attention page size, causing severe capacity overestimation documented in vLLM issue \#37121 (concrete magnitudes in \S\ref{sec:method:avpt}). Alternatively, SGLang implements a static dual pool where operators fix a \texttt{mamba\_full\_memory\_ratio} parameter at engine initialization \cite{zheng2024sglang}. This rigid partition cannot rebalance without a full process restart.

Neither architectural approach handles workload shifts dynamically. Static defaults perform poorly when prompt distributions change during execution. Testing the same workload mix records 1221 Out-of-Memory (OOM) events at a 0.9 static ratio compared to 552 OOMs at a 0.5 static ratio. No existing system provides both asymmetric cache types and runtime adaptivity.

In this paper, we introduce Asymmetric Virtual Memory Paging (AVMP), a paged memory allocator designed specifically for hybrid architectures. AVMP provisions physically asymmetric backing stores while enabling dynamic capacity rebalancing at runtime. Our dynamic rebalancing records a 13.3$\times$ goodput improvement on \texttt{uniform\_short} workloads compared to the best static baseline.

\subsection{Contributions}
\label{sec:intro:contributions}

We make the following specific contributions:

\begin{itemize}[leftmargin=1.4em, itemsep=2pt, topsep=4pt]
\item \textbf{Asymmetric Virtual Page Table Abstraction (\S\ref{sec:method}):} We design a unified virtual handle space spanning two physically heterogeneous backing stores (KV pages and SSM blocks). AVMP extends the GPU virtual memory approach pioneered for KV caches \cite{xu2024vtensor} to two asymmetric pool types. We show that this abstraction introduces no measurable effect on simulated capacity outcomes, with byte-identical OOM counts against a static dual-pool baseline (552 OOMs).

\item \textbf{Dynamic Pool Rebalancing Mechanism (\S\ref{sec:method}):} We implement a \texttt{CapacityError}-triggered migration mechanism that transfers batched memory capacity between pools. We enforce determinism using a logical operation counter rather than wall-clock time, ensuring per-cell byte-identical reproducibility across all experimental reruns.

\item \textbf{Empirical Validation on Hybrid Architectures (\S\ref{sec:eval}):} We execute a 270-cell sweep evaluating 5 allocator variants, 3 synthetic workloads, 3 pool sizes, and 3 random seeds. The dynamic AVMP allocator records a 7.6\% OOM reduction (510 versus 552) against the best static baseline. AVMP records 13.3$\times$ higher goodput on \texttt{uniform\_short}, 2.39$\times$ on \texttt{mixed\_long}, and 1.83$\times$ on \texttt{agentic\_burst} traffic patterns. A parameter sensitivity analysis proves \texttt{migration\_batch\_size} acts as the dominant configuration axis, while a Stage 2 threshold sweep returned a strict null result.

\item \textbf{Open-Source Prototype:} We provide a reference Python implementation and synthetic workload harness with committed sweep artifacts at \url{https://github.com/codepawl/cachepawl}.
\end{itemize}
